% SOURCE_AUTHORITY: sga5_fr_workpass.tex, lines 14082--14116 (LF-normalized unit).
% SOURCE_SHA256: 791F4EFFC5E02832D5D77ED03518C8156D6F07E4C8238B03545DB93D883FBB28
\subsubsection*{d) O caso em que o feixe $F$ não é ramificado nos pontos fixos do endomorfismo $f$.}

Neste caso, o termo local $\tau_x(f,u, f', G, \varphi)$ é dado pela fórmula:

\begin{equation}\tag{6.4.6}
\tau_x(f,u, f', G, \varphi) = \nu_x(f) \, \Tr_A(u_x) \quad ,
\end{equation}
e, por conseguinte, depende apenas do par $(f,u)$.

Com efeito, neste caso, tem-se:

\begin{equation}\tag{6.4.7}
S'_{x,\ell^n}(f, f', G, \varphi)(g) = \begin{cases}
0 & \text{se existe } x' \in X' \text{ tal que } (g^{-1}f^{\prime})(x') = x' \text{ ou}\\
& \text{equivalentemente, se a classe de $\varphi$-}\\
& \text{conjugação de $g$ coincide com}\\
& \text{a classe } \mathrm{NW}_{f}^{G,\varphi}(x) \text{ (cf. 3); }\\
(\nu_x(f)-1) & \text{caso contrário.}
\end{cases}
\end{equation}

Em particular, resulta que:

% SOURCE_ANOMALY_PRESERVED: (6.4.7) uses S' but (6.4.8) drops the prime and introduces a semicolon after f'; frozen authority and printed p.437 agree.
\begin{equation}\tag{6.4.8}
S_{x,\ell^n}(f, f'; G, \varphi)(g) = 0 \quad \text{se } g \notin \mathrm{NW}_{f}^{G,\varphi}(x) \quad .
\end{equation}

Obtém-se, portanto:

\begin{equation}\tag{6.4.9}
\begin{aligned}
\tau_x(f,u, f', G, \varphi) &= \nu_x(f) \Tr_A(g^{-1}u_M^\vee) = \nu_x(f) \Tr_A(g^{-1}u')_{x'} \\
&= \nu_x(f) \Tr_A(u_x) \quad .
\end{aligned}
\end{equation}
